% GENERATED FILE. Edit canonical lessons, then run npm run generate:books.
% canonical-source-hash: fnv1a64:4929d55380341bad
% canonical-lessons: GE-C142-zukunft, GE-C142-naechste-woche, GE-C142-uebermorgen, GE-C142-naechstes-jahr, GE-C142-vorhaben

\chapter{The Future, Next Week, the Day after Tomorrow, Next Year, To Plan}
\label{ch:ge-the-future-next-week-the-day-after-tomorrow-next-year-to-plan}
\hlchaptermodality{142}

\begin{chapteropening}
\textbf{By the end of this chapter:} \emph{I can say what will happen, and when: next week, the day after tomorrow, next year.}

Complete the last lesson of chapter 142: i can say what will happen, and when: next week, the day after tomorrow, next year.
\end{chapteropening}

\section[die Zukunft]{die Zukunft — the future}

\label{lesson:GE-C142-zukunft}



\begin{quote}
Before the new one: say the German for to repair, then the German for to start.
\end{quote}

\subsection*{You'll want to know: die Zukunft}
\textbf{die Zukunft} — \textquotedblleft{}the future\textquotedblright{}. German builds the future with \textbf{werden} and an infinitive at the very end: \textbf{Ich werde Deutsch lernen.} — \textquotedblleft{}I will learn German.\textquotedblright{}

\begin{grammarlens}[title={What you've built}]
One more word for this chapter.
\end{grammarlens}

\subsection*{Your turn}
\begin{itemize}
\raggedright
  \item \emph{Say it:} \emph{die Zukunft}
  \item \emph{Say it:} \emph{die Zukunft}, once more
  \item \emph{Say it:} say \emph{die Zukunft}
  \item \emph{Recall:} say the German for to repair, then the German for to start, then say \emph{die Zukunft} again
\end{itemize}

\begin{tcolorbox}[breakable,colback=teal!4,colframe=teal!35!black,title={Before you move on}]
What does die Zukunft mean? (\textquotedblleft{}The future\textquotedblright{}.) Say it once more.
\end{tcolorbox}
\section[nächste Woche]{nächste Woche — next week}

\label{lesson:GE-C142-naechste-woche}



\begin{quote}
Before the new one: say the German for to start, then the German for the future.
\end{quote}

\subsection*{You'll want to know: nächste Woche}
\textbf{nächste Woche} — \textquotedblleft{}next week\textquotedblright{}. Put the time first and the verb stays second: \textbf{Nächste Woche werde ich Deutsch lernen.} \emph{Werde} comes before \emph{ich}, and \emph{lernen} waits at the end.

\begin{grammarlens}[title={What you've built}]
One more word for this chapter.
\end{grammarlens}

\subsection*{Your turn}
\begin{itemize}
\raggedright
  \item \emph{Say it:} \emph{nächste Woche}
  \item \emph{Say it:} \emph{nächste Woche}, once more
  \item \emph{Say it:} say \emph{nächste Woche}
  \item \emph{Recall:} say the German for to start, then the German for the future, then say \emph{nächste Woche} again
\end{itemize}

\begin{tcolorbox}[breakable,colback=teal!4,colframe=teal!35!black,title={Before you move on}]
What does nächste Woche mean? (\textquotedblleft{}Next week\textquotedblright{}.) Say it once more.
\end{tcolorbox}
\section[übermorgen]{übermorgen — the day after tomorrow}

\label{lesson:GE-C142-uebermorgen}



\begin{quote}
Before the new one: say the German for the future, then the German for next week.
\end{quote}

\subsection*{You'll want to know: übermorgen}
\textbf{übermorgen} — \textquotedblleft{}the day after tomorrow\textquotedblright{}. When a time word already points ahead, spoken German keeps the present tense: \textbf{Übermorgen fahre ich.} — \textquotedblleft{}I am going the day after tomorrow.\textquotedblright{}

\begin{grammarlens}[title={What you've built}]
One more word for this chapter.
\end{grammarlens}

\subsection*{Your turn}
\begin{itemize}
\raggedright
  \item \emph{Say it:} \emph{übermorgen}
  \item \emph{Say it:} \emph{übermorgen}, once more
  \item \emph{Say it:} say \emph{übermorgen}
  \item \emph{Recall:} say the German for the future, then the German for next week, then say \emph{übermorgen} again
\end{itemize}

\begin{tcolorbox}[breakable,colback=teal!4,colframe=teal!35!black,title={Before you move on}]
What does übermorgen mean? (\textquotedblleft{}The day after tomorrow\textquotedblright{}.) Say it once more.
\end{tcolorbox}
\section[nächstes Jahr]{nächstes Jahr — next year}

\label{lesson:GE-C142-naechstes-jahr}



\begin{quote}
Before the new one: say the German for next week, then the German for the day after tomorrow.
\end{quote}

\subsection*{You'll want to know: nächstes Jahr}
\textbf{nächstes Jahr} — \textquotedblleft{}next year\textquotedblright{}. \emph{Werden} changes with the person: \emph{ich werde}, \emph{du wirst}, \emph{er wird}. \textbf{Nächstes Jahr wirst du Deutsch lernen.}

\begin{grammarlens}[title={What you've built}]
One more word for this chapter.
\end{grammarlens}

\subsection*{Your turn}
\begin{itemize}
\raggedright
  \item \emph{Say it:} \emph{nächstes Jahr}
  \item \emph{Say it:} \emph{nächstes Jahr}, once more
  \item \emph{Say it:} say \emph{nächstes Jahr}
  \item \emph{Recall:} say the German for next week, then the German for the day after tomorrow, then say \emph{nächstes Jahr} again
\end{itemize}

\begin{tcolorbox}[breakable,colback=teal!4,colframe=teal!35!black,title={Before you move on}]
What does nächstes Jahr mean? (\textquotedblleft{}Next year\textquotedblright{}.) Say it once more.
\end{tcolorbox}
\section[vorhaben]{vorhaben — to plan, to intend}

\label{lesson:GE-C142-vorhaben}



\begin{quote}
Before the new one: say the German for the day after tomorrow, then the German for next year.
\end{quote}

\subsection*{You'll want to know: vorhaben}
\textbf{vorhaben} — \textquotedblleft{}to plan, to intend\textquotedblright{}. It splits like the other separable verbs: \textbf{Was hast du vor?} — \textquotedblleft{}What are you planning?\textquotedblright{}

\begin{grammarlens}[title={What you've built}]
One more word for this chapter.
\end{grammarlens}

\subsection*{Your turn}
\begin{itemize}
\raggedright
  \item \emph{Say it:} \emph{vorhaben}
  \item \emph{Say it:} \emph{vorhaben}, once more
  \item \emph{Say it:} say \emph{vorhaben}
  \item \emph{Recall:} say the German for the day after tomorrow, then the German for next year, then say \emph{vorhaben} again
\end{itemize}

\begin{tcolorbox}[breakable,colback=teal!4,colframe=teal!35!black,title={Before you move on}]
What does vorhaben mean? (\textquotedblleft{}To plan, to intend\textquotedblright{}.) Say it once more.
\end{tcolorbox}
